% SOURCE_STATEMENT: "How is the apparent paradox between topographic Rossby waves and inertial waves resolved: does the series whose zeroth-order term is the topographic Rossby wave diverge locally, or does it coexist with the inertial wave?"
% Preparation and final self-audit: CONTRIBUTING.md prompts C and D; source/audit artifacts retained in the local physics_statements_2026_10_11 batch.
% !TeX program = lualatex
\documentclass[11pt]{article}
\usepackage[margin=25mm]{geometry}
\usepackage{fontspec}
\setmainfont{Latin Modern Roman}
\usepackage{amsmath,amssymb,mathtools}
\usepackage{unicode-math}
\setmathfont{Latin Modern Math}
\usepackage{xurl,hyperref}
\hypersetup{colorlinks=true,urlcolor=blue}
\setcounter{secnumdepth}{0}
\setlength{\parindent}{0pt}
\setlength{\parskip}{5pt}
\setlength{\emergencystretch}{3em}
\newcommand{\sref}[2]{\hyperlink{source:#1:#2}{[S#2]}}
\newcommand{\cataloguescope}[1]{\paragraph{Recorded target.}#1}
\newcommand{\R}{\mathbb R}
\newcommand{\C}{\mathbb C}
\newcommand{\N}{\mathbb N}
\newcommand{\Z}{\mathbb Z}
\newcommand{\Q}{\mathbb Q}
\newcommand{\F}{\mathbb F}
\newcommand{\eps}{\varepsilon}
\begin{document}
\section{Topographic Rossby Waves and Inertial-Wave Coexistence}\label{statement}

\subsection{Definitions and mathematical statement}

\paragraph{Shared notation.}$\mathbb N=\{1,2,\ldots\}$; the standard number fields have their usual meanings. All additional conventions are stated locally.


\paragraph{Definitions and conventions.} A rotating fluid supports inertial waves due to Coriolis restoring forces. A sloping boundary or varying depth gives an effective topographic potential-vorticity gradient and supports topographic Rossby waves. The report compares these waves in a rectangular rotating basin with a sloping boundary and examines an asymptotic series whose leading term is a topographic Rossby wave.

\paragraph{Source scope and qualifications.} The basin geometry, boundary slope, and asymptotic expansion of the contribution are part of the context. The source does not ask whether every Rossby wave in every geometry diverges; the two alternatives concern the particular paradox described there.

\paragraph{Catalogue mathematical target.}
\paragraph{Catalogue wording retained for traceability.}
How is the apparent paradox between topographic Rossby waves and inertial waves resolved: does the series whose zeroth-order term is the topographic Rossby wave diverge locally, or does it coexist with the inertial wave?

\paragraph{Saved literature context.} The supplied catalogue's assessment dated 2026-08-21 records: The report presents this as an open problem, conjecture, or unresolved research direction; targeted literature checking found no authoritative resolution of the extracted claim as of 2026-08-21. The assessment is retained as dated provenance; it is not a new comprehensive literature or open-status certification.
\subsection{Short English statement}
Explain whether the topographic-Rossby expansion breaks down locally or describes waves coexisting with the exact inertial waves.
\subsection{Sources}
\begin{itemize}
\item[\hypertarget{source:statement:1}{S1}] ULAM.AI and the UnsolvedMath Contributors, UnsolvedMath, catalogue entry. Catalogue curation under CC BY 4.0; definitions and English exposition added locally.
\url{https://huggingface.co/datasets/ulamai/UnsolvedMath}.
\item[\hypertarget{source:statement:2}{S2}] Mathematical Theory and Modelling in Atmosphere-Ocean Science (2006), Oberwolfach Reports, DOI 10.4171/owr/2006/39, pp. 2307--2398; official report source and mathematical context.
\url{https://doi.org/10.4171/owr/2006/39}.
\item[\hypertarget{source:statement:3}{S3}] Official EMS Press PDF of the same Oberwolfach report; PDF page 75.
\url{https://ems.press/content/serial-article-files/46068}.
\end{itemize}
\label{end:statement}
\end{document}
